\chapter{Conclusões e Trabalhos Futuros}
\label{chap:conclusoes-e-trabalhos-futuros}


A partir desses resultados fica evidente como os modelos de Inteligência Artificial merecem sua devida atenção, pois conseguiram performar tão bem quanto modelos tradicionais consolidados no mercado. Contudo mesmo diante de um desempenho formidável desses modelos "novos" é de extrema relevância ressaltar que essas novas metodologias não vieram para substituir os chamados modelos tradicionais, mas sim para complementá-las.

Os testes realizados durante esse trabalho podem ser aplicados em contextos reais de problemas econômicos empresarias e governamentais. E nesses ambientes, onde os dados estão presentes em grande volume, a metodologia deixa de ser encontrar o melhor modelo, mas sim testar com todos os presentes, e entender se as respostas apontam para o mesmo caminho. Só com essa informação em mãos os tomadores de decisão poderão decidir entre comprar e vende, investir ou poupar, consumir ou adiar. 

Essas decisões que moldam o pensamento econômico quase nunca são estritamente matemáticas, não se trata de uma fórmula mágica, muitas vezes dependem de um feeling, um estalo, uma decisão final baseada em conhecimento tácito sobre o negócio. Esse tipo de decisão não pode ser determinado por relações lineares de modelos tradicionais. Muitas vezes não por uma decisão inteligente, mas por uma falta de capacidade, essas decisões são descartadas por um modelo. Contudo conforme apresentado, os modelos de IA atacam justamente esse ponto, a relevância, o fator não-linear nas relações econômicas. 

Existe uma ironia ao refletir sobre esse tema, as relações humanas são tão complexas, que precisamos de uma máquina, que não entende essas relações, para nos auxiliar. Mas não se trata de adivinhação, não é mágica, são as relações humanas históricas do passado que fundamentam as predições futuras. Muito se discute se as máquinas um dia serão inteligentes e pensarão como humanos, porém essa discussão não é relevante para esse problema. O que precisamos realmente discutir é, as máquinas conseguem extrair de um conjunto de dados, não a tendência mecânica e linear de uma economias, mas sim o comportamento humano. 

Tendo esse momento de reflexão, é fato também relatar o desempenho dos modelos, é notável que a econometria tem um trabalho relevante anterior ao modelo, a seleção, coleta, tratamento de variáveis é o procedimento mais trabalhoso da aplicação aos modelos. E em sua totalidade é responsável pelo desempenho do modelo e as conclusões que um pesquisador terá sobre o tema. Durante essa pesquisa a proposta inicial era realizar com mais variáveis, mas durante os tratamentos e decisões de cortes, e a própria arquitetura limitante de alguns modelos, retinge o uso de variáveis.

Além é claro dos problemas enfrentados no que tange a disponibilidade de dados e continuidade das séries, o próprio PIB anual que precisou ser reconstruído com base nos valores trimestrais do IBGE, a falha de continuidade na série do desemprego devido a mudança metodológica, sem falar da disponibilidade reduzida de observações para estimar com confiança uma determinada variável alvo. 

Nesse sentido poderíamos testar os modelos com dados fictícios, porém isso primeiro reduziria a perspectiva econômica do trabalho, que acaba sendo parte fundamental, tendo em vista a área de estudo a a finalidade desse trabalho. Porém a testagem laboratorial, sem limitações em condições controladas, tem muito menos a oferecer no quesito realidade. Esse processo de tratamento e escolha de variáveis, a avaliação do comportamento dos modelo quando um problema real é adicionado, seja um dado faltante, uma medição não confiável, uma quebra metodológica, o modelo consegue estimar o alvo com precisão, ou a sensibilidade ao erro é um fator crucial.

Porém não só de problemas vive esse trabalho, mas também de solucioná-los, ou pelo mínimo apontar uma possibilidade no emaranhado de respostas existentes. E o problema que teve uma das respostas apontadas por esse trabalho é como desenvolver um modelo preditivo baseado em IA. Essa explicação não é simples, na realidade ela passou por uma caminho trilhado em 4 objetivos: Desde a caracterização dos modelos que ocorreu em grande parte na revisão de literatura, mas também na apresentação de nossos resultados; A seleção das variáveis tendo em vista as limitações de coleta e período; A construção do modelo em si que envolveu a testagem de diversas variações até encontrar a estrutura mais adequadas aos dados utilizados; Por fim, é claro a apresentação do desempenho de cada modelo, comparativamente, estabelecendo não um pior, ou melhor absoluto, mas um desempenho estrito aos nossos dados e nossas variáveis, que podem futuramente serem complementados e melhorados.

E é a partir dessa complementação e melhoria que esse trabalho pode continuar, o espaço que se abre para uma pesquisa futura é a interpretação de modelos de IA, um das grandes restrições das redes neurais especificamente são as interpretações dos pesos dos modelos. Em aplicações onde a transparência é uma fator relevante como em decisões governamentais, ou decisões bancárias por exemplo, a interpretação é tão importante quanto o acerto. Nesse sentido a avaliação de interpretabilidade é fundamental.

Outro ramo de pesquisa futuro é o campo das IAs generativas, essas tem um espaço para pesquisa muito vasto e ainda muito inexplorado no âmbito da pesquisa acadêmica. Enquanto as IAs preditiva utilizam os dados históricos para classificar, prever ou tomar decisões baseadas em regras. As IAs generativas tem uma capacidade de sintetizar o conhecimento e gerar "novos conteúdos" baseados no aprendizado durante o treinamento. Isso possui um amplo campo de estudo futuro, em que a combinação dos modelos tradicionais, modelos de IAs preditivas e generativas pode ser o trabalho do economista do futuro.

De toda forma a importância do estudo acadêmico das máquinas econômicas é fundamental, o uso da IA combinado com métodos tradicionais e com o olhar do economista comandando os parâmetros, escolhendo os dados e ajustando os controles é o futuro. E respondendo a pergunta inicial desse trabalho, e parafraseando Keynes. Para onde vamos? não sabemos, mas estamos fazendo algo, para continuar indo para lá.





